% Class-compatibility preamble retained for a reader excerpt of
%   Dave Witte Morris, "A Cohomological Proof that Real Representations of
%   Semisimple Lie Algebras Have Q-Forms", SIGMA 11 (2015) 034,
%   DOI 10.3842/SIGMA.2015.034, CC BY 4.0.
%
% The source pins the publisher class sigma, which is not part of the delivery
% baseline, so this file supplies the amsart baseline and then copies the author's
% own declaration block mechanically and unchanged:
%   - pinned-original-source.tex lines 11-31: the author's DeclareMathOperator
%     list, the equation numbering and the author's newtheorem declarations;
%   - pinned-original-source.tex lines 35-43, which the author writes immediately
%     after \begin{document} in the publisher layout: the author's own \Res,
%     \ResOne and \Lie macros and \allowdisplaybreaks.
% The publisher header commands of the sigma class (\FirstPageHeading,
% \ArticleName and the like) and the macro \PaperNumber they require are not
% reproduced, because this excerpt carries neither the journal front matter nor
% anything else that uses them.
% No mathematical macro, statement, symbol or assumption is introduced, changed
% or removed by this file.  The author's full DeclareMathOperator list of pinned
% source lines 11-19 is carried, including GL, SL, SO and SU: these four are not
% supplied by the delivery baseline, so leaving them out makes the author's own
% notation (for instance GL_n(Q) in the corollary on lattices) undefined.
\documentclass[10pt,a4paper]{amsart}
\usepackage{amssymb,amsmath,amsthm}
\usepackage{enumerate}
\usepackage{hyperref}

\DeclareMathOperator{\Cliff}{Clif\/f}
\DeclareMathOperator{\Gal}{Gal}
\DeclareMathOperator{\GL}{GL}
\DeclareMathOperator{\Inn}{Inn}
\DeclareMathOperator{\Mat}{Mat}
\DeclareMathOperator{\Out}{Out}
\DeclareMathOperator{\SL}{SL}
\DeclareMathOperator{\SO}{SO}
\DeclareMathOperator{\SU}{\mathbf{SU}}
\numberwithin{equation}{section}

\newtheorem{prop}{Proposition}[section]
\newtheorem{cor}[prop]{Corollary}
\newtheorem{lem}[prop]{Lemma}
{\theoremstyle{definition}
\newtheorem{defn}[prop]{Definition}
\newtheorem{notation}[prop]{Notation}
\newtheorem{rem}[prop]{Remark}
\newtheorem{ques}[prop]{Question}
\newtheorem{eg}[prop]{Example}
\newtheorem{erratum}[prop]{Erratum}}

\newcommand{\Res}[1]{\mathop{\mathbf{R}_{#1}}}
\newcommand{\ResOne}[1]{\mathop{\mathbf{R}^{(1)}_{#1}}}
\newcommand{\Lie}[1]{\mathfrak{\lowercase{#1}}}
\allowdisplaybreaks
\usepackage{cleveref}
\crefname{originallemma}{original lemma}{original lemmas}
\Crefname{originallemma}{Original lemma}{Original lemmas}
\makeatletter
\newlabel{RunivIff}{{2.1}{source}{}{sourcepointer-1}{}}
\newlabel{RunivIff@cref}{{[originallemma][0][]2.1}{[1][0][]source}{}{}{}}
\makeatother
\crefname{originaldefinition}{original definition}{original definitions}
\Crefname{originaldefinition}{Original definition}{Original definitions}
\makeatletter
\newlabel{RunivIff-DivAlg}{{3}{source}{}{sourcepointer-2}{}}
\newlabel{RunivIff-DivAlg@cref}{{[originaldefinition][0][]3}{[1][0][]source}{}{}{}}
\makeatother
\crefname{originallemma}{original lemma}{original lemmas}
\Crefname{originallemma}{Original lemma}{Original lemmas}
\makeatletter
\newlabel{SameFixed}{{2.3}{source}{}{sourcepointer-3}{}}
\newlabel{SameFixed@cref}{{[originallemma][0][]2.3}{[1][0][]source}{}{}{}}
\makeatother
\title{R-universality of the even special orthogonal algebra in terms of the even Clifford algebra and the determinant sign}
\author{Dave Witte Morris}
\date{doi:10.3842/SIGMA.2015.034; CC BY 4.0}
\begin{document}
\maketitle
\noindent\textit{Editorial note.} All delivered mathematical text is the author's own, copied from the pinned publisher source sources/algebra-fl02-a-cohomological-proof-that-real-re/pinned-original-source.tex; no statement or proof step was rewritten, reordered or removed.\par\medskip
\noindent\textit{Editorial note.} Preamble provenance: the source pins the publisher class sigma, which is not part of the delivery baseline, so this card uses sources/algebra-fl02-a-cohomological-proof-that-real-re/retained-preamble.tex, which supplies the amsart baseline and then carries the author's own declarations from the pinned source (its header records the exact line ranges). The author's full DeclareMathOperator list is carried, including GL, SL, SO and SU, because the delivery baseline does not provide those four operator names. No macro was invented and no mathematical symbol was changed.\par\medskip
\noindent\textit{Editorial note.} Bibliography provenance: the citations of the counted proof are resolved with sources/algebra-fl02-a-cohomological-proof-that-real-re/bibliography.tex, an excerpt of the author's own thebibliography environment (pinned source lines 954-1016) copied character for character with all twelve entries kept; no entry was rewritten, reordered, merged, split or invented and no citation key was changed.\par\medskip
\noindent\textit{Editorial note.} External references are navigation pointers only; every deliberately kept step of the author's argument stays inside the counted spans.\par\medskip
\noindent\textit{Editorial note.} Selection provenance: this card was opened under the file-based control contract (CONTROL.md), with the mechanical pre-card gate run on its statement and proof interval before the card was written.\par\medskip
\noindent\textit{Allowed standard prerequisites (editorial).}\begin{itemize}
\item The Tits algebras attached to the three non-zero elements of the character group for type D\_n, and the identification of the first of them with the even Clifford algebra. Used inside the counted proof at source lines 684-690 and 708-711, attributed by the author to section 27B, type D\_n, page 379 of the cited book KnusEtAl-BookOfInvols. Supplies the Brauer classes whose base-change behaviour decides R-universality.
\item The division-algebra description of R-universality and the meaning of being an inner form for these orthogonal algebras. The author's lemma labelled RunivIff, invoked in the counted proof at source lines 710-711, together with the standard description of inner forms of type D\_n. Turns the computation of the Tits algebras into the two alternatives of the statement.
\end{itemize}\medskip
\ifdefined\hypertarget\hypertarget{sourcepointer-1}{}\fi
\noindent\textit{Source reference.} Original lemma 2.1: sources/algebra-fl02-a-cohomological-proof-that-real-re/pinned-original-source.tex, source lines 416--416. Navigation pointer only: the counted proof invokes the author's equivalence for R-universality; that lemma has its own card in this lane and its content is declared as a named standard prerequisite.\par\medskip
\ifdefined\hypertarget\hypertarget{sourcepointer-2}{}\fi
\noindent\textit{Source reference.} Original definition 3: sources/algebra-fl02-a-cohomological-proof-that-real-re/pinned-original-source.tex, source lines 428--428. Navigation pointer only: the labelled item of that equivalence used by the counted proof.\par\medskip
\ifdefined\hypertarget\hypertarget{sourcepointer-3}{}\fi
\noindent\textit{Source reference.} Original lemma 2.3: sources/algebra-fl02-a-cohomological-proof-that-real-re/pinned-original-source.tex, source lines 471--471. Navigation pointer only: the author's lemma on the minimal inner extension; it has its own card in this lane and its content is declared as a named standard prerequisite.\par\medskip
\begin{prop}
\label{so2k}
Let $\Lie G = \Lie{SO}_{2k}(B; \mathbb{Q})$, with $k \ge 3$, for some symmetric $B \in \GL_{2k}(\mathbb{Q})$.
Then~$\Lie G$ is $\mathbb{R}$-universal if and only if either
\renewcommand{\labelenumi}{{\rm (\theenumi)}}
\begin{enumerate}\itemsep=0pt
\item $\Cliff^0_\mathbb{Q}(B)$ is split, and $(-1)^k \det B$ is either negative or a~square in~$\mathbb{Q}$, or
\item $\Cliff^0_\mathbb{Q}(B)$ does not split over~$\mathbb{R}$, and $(-1)^k \det B$ is a~square in~$\mathbb{Q}$.
\end{enumerate}
\end{prop}

\begin{proof}
As in~\cite[\S~27B, type~$D_n$, p.~379]{KnusEtAl-BookOfInvols}, let $\lambda$, $\lambda_+$, $\lambda_-$ be dominant weights
that represent the three nonzero elements of~$Z^*$, with~$\lambda$ being the highest weight of the natural
representation of~$\Lie G$ on $\mathbb{C}^{2k}$.
Then $D_{\Lie G, \mathbb{Q}}(\lambda) = \mathbb{Q}$ is trivial~\cite[p.~379]{KnusEtAl-BookOfInvols}.

Suppose, f\/irst, that $\Lie G$ is an inner form, which means that $(-1)^k \det B$ is a~square in~$\mathbb{Q}$.
Then $\Cliff^0_\mathbb{Q}(B)$ is a~direct sum of two algebras $C^+$ and~$C^-$ that are Brauer equivalent to the full
Clif\/ford algebra~\cite[Theorem~2.5(3), p.~110]{Lam-IntroQuadForms} (and are therefore Brauer equivalent to a~quaternion
algebra).
Furthermore, $D_{\Lie G, \mathbb{Q}}(\lambda_{\pm})$ is Brauer equivalent
to~$C^\pm$~\cite[p.~379]{KnusEtAl-BookOfInvols}.
Therefore, Lemma~\ref{RunivIff}(\ref{RunivIff-DivAlg}) shows that an inner form $\Lie G$ is:
\begin{itemize}\itemsep=0pt
\item automatically $\mathbb{R}$-universal, when $\Cliff^0_\mathbb{Q}(B)$ does not split over~$\mathbb{R}$, but
\item $\mathbb{R}$-universal if and only if $\Cliff^0_\mathbb{Q}(B)$ is split, when $\Cliff^0_\mathbb{Q}(B)$ splits
over~$\mathbb{R}$.
\end{itemize}

Assume, now, that $\Lie G$ is an outer form.
In order for $\Lie G$ to be $\mathbb{R}$-universal, $\Lie G$ must remain outer over~$\mathbb{R}$ (see Lemma~\ref{SameFixed}),
which means $(-1)^k \det B < 0$.
Let~$L$ be the (unique, imaginary) quadratic extension of~$\mathbb{Q}$ over which $\Lie G$ becomes inner.
Then $D_{\Lie G, \mathbb{Q}}(\lambda_{\pm})$ is Brauer equivalent to
$\Cliff^0_\mathbb{Q}(B)$~\cite[p.~379]{KnusEtAl-BookOfInvols}, which is central simple over~$L$.
Since~$B$ splits over~$\mathbb{R}$ (recall that the quadratic extension~$L$ is
imaginary), Lemma~\ref{RunivIff}(\ref{RunivIff-DivAlg}) implies that $\Lie G$ is $\mathbb{R}$-universal if and only if
$\Cliff^0_\mathbb{Q}(B)$ is split.
\end{proof}

%% Bibliography excerpt retained for reader excerpts of
%   Dave Witte Morris, "A Cohomological Proof that Real Representations of Semisimple Lie Algebras Have Q-Forms",
%   SIGMA 11 (2015) 034, DOI 10.3842/SIGMA.2015.034.
%%
%% Every entry below is copied character for character from the author's own
% thebibliography environment of the pinned publisher source (lines 954-1016).
%% No entry was rewritten, reordered, merged, split or invented, and no citation
% key was changed.  All 12 entries are kept so that any citation of a delivered
%% excerpt resolves without guessing an entry.
\begin{thebibliography}{99}

\footnotesize\itemsep=0pt

\bibitem{DemazureGrothendieck-SGA3}
Demazure M., Grothendieck A., Sch\'emas en groups.~III. Structure des sch\'emas
  en groupes reductifs, S\'eminaire de g\'eom\'etrie alg\'ebrique du {B}ois
  {M}arie, 1962/1964, available at
  \url{http://library.msri.org/books/sga/sga/pdf/sga3-3.pdf}.

\bibitem{Harder-BerichtNeuere}
Harder G., Bericht \"uber neuere {R}esultate der {G}aloiskohomologie
  halbeinfacher {G}ruppen, \textit{Jber. Deutsch. Math.-Verein.} \textbf{70}
  (1968), 182--216, available at
  \url{http://www.digizeitschriften.de/en/dms/img/?PPN=PPN37721857X_0070&DMDID=dmdlog17}.

\bibitem{Kneser-LectGalCohoClassical}
Kneser M., Lectures on {G}alois cohomology of classical groups, \textit{Tata
  Institute of Fundamental Research Lectures on Mathematics}, Vol.~47, Tata
  Institute of Fundamental Research, Bombay, 1969, available at
  \url{http://www.math.tifr.res.in/~publ/ln/tifr47.pdf}.

\bibitem{KnusEtAl-BookOfInvols}
Knus M.A., Merkurjev A., Rost M., Tignol J.P., The book of involutions,
  \textit{American Mathematical Society Colloquium Publications}, Vol.~44,
  Amer. Math. Soc., Providence, RI, 1998.

\bibitem{Lam-IntroQuadForms}
Lam T.Y., Introduction to quadratic forms over f\/ields, \textit{Graduate Studies
  in Mathematics}, Vol.~67, Amer. Math. Soc., Providence, RI, 2005.

\bibitem{Morris-RRepsQForms}
Morris D., Real representations of semisimple {L}ie algebras have {${\mathbb
  Q}$}-forms, in Algebraic Groups and Arithmetic, Tata Institute of Fundamental
  Research, Mumbai, 2004, 469--490.

\bibitem{PlatonovRapinchukBook}
Platonov V., Rapinchuk A., Algebraic groups and number theory, \textit{Pure and
  Applied Mathematics}, Vol.~139, Academic Press, Inc., Boston, MA, 1994.

\bibitem{PrasadRapinchuk-Prescribed}
Prasad G., Rapinchuk A., On the existence of isotropic forms of semi-simple
  algebraic groups over number f\/ields with prescribed local behavior,
  \href{http://dx.doi.org/10.1016/j.aim.2006.01.001}{\textit{Adv. Math.}} \textbf{207} (2006), 646--660.

\bibitem{Raghunathan-ArithLattsSSGrps}
Raghunathan M.S., Arithmetic lattices in semisimple groups, \href{http://dx.doi.org/10.1007/BF02967980}{\textit{Proc.
  Indian Acad. Sci. Math. Sci.}} \textbf{91} (1982), 133--138.

\bibitem{Tits-Classification}
Tits J., Classif\/ication of algebraic semisimple groups, in Algebraic {G}roups
  and {D}iscontinuous {S}ubgroups ({P}roc. {S}ympos. {P}ure {M}ath., {B}oulder,
  {C}olo., 1965), Amer. Math. Soc., Providence, R.I., 1966, 33--62.

\bibitem{Tits-RepLinIrred}
Tits J., Repr\'esentations lin\'eaires irr\'eductibles d'un groupe r\'eductif
  sur un corps quelconque, \href{http://dx.doi.org/10.1515/crll.1971.247.196}{\textit{J.~Reine Angew. Math.}} \textbf{247} (1971),
  196--220.

\bibitem{Tits-StronglyInner}
Tits J., Strongly inner anisotropic forms of simple algebraic groups,
  \href{http://dx.doi.org/10.1016/0021-8693(90)90201-X}{\textit{J.~Algebra}} \textbf{131} (1990), 648--677.


\end{thebibliography}
\end{document}
